\section{Purpose}

\textbf{[STIPULATION]} This document develops the Dirac equation in FFGFT
completely: the Clifford algebra as fundamental structure, the transition from the
standard Dirac to the T0 formulation, mass elimination (physics formulable without
mass parameters), spin as topological property, and the modified Dirac equation
from the time-field Lagrangian (A142). The document supplements A075 (field theory,
Schrödinger limit) with the complete relativistic structure.

\section{The Clifford Algebra as Fundamental Basis}

\textbf{[B]} The standard Dirac equation $(i\gamma^\mu\partial_\mu - m)\psi=0$
uses $4\times4$ $\gamma$-matrices as a representation. In FFGFT the fundamental
structure is the Clifford algebra:
\[
  \mathbf e_\mu\mathbf e_\nu + \mathbf e_\nu\mathbf e_\mu = 2g_{\mu\nu}.
\]
The $\gamma$-matrices are one representation -- they are not eliminated but traced
back to their geometric origin. That is the difference: in the SM $\gamma$-matrices
are the foundation; in FFGFT it is the Clifford algebra, with $\gamma$-matrices
being only one representation.

\textbf{[B]} \textbf{What changes in T0?} The standard Dirac equation has $m$ as a
fixed parameter. In T0 the mass $m(x,t)$ is a dynamic field:
\[
  \bigl[i\gamma^\mu(\partial_\mu + \Gamma_\mu^{(T)}) - m(x,t)\bigr]\psi = 0,
\]
with time-field connection $\Gamma_\mu^{(T)}=\partial_\mu m/m^2$ (A142). In the
limit $\partial_\mu m\to0$ (constant mass) this reduces to the standard Dirac
equation.

\section{Spin as Topological Property}

\textbf{[B]} In FFGFT spin is not an intrinsic property but a topological one:
spin-$\tfrac12$ corresponds to the torus winding $(n_\theta,n_\phi)=(1,1)$. The
720° symmetry (two rotations return a spinor to its initial state) follows directly
from the topology of the torus -- it need not be postulated.

\textbf{[B]} A relativistic spinor state in the torus representation:
\[
  \Phi(x,t) = \rho(x,t)\,e^{i\theta(x,t)}\cdot\chi(\sigma),
\]
with spin component $\chi(\sigma)$ from the geometric knot modes (A075). The quantum
number $j=l\pm\tfrac12$ for fermions follows from the winding structure.

\section{Mass Elimination: Physics Without Mass Parameters}

\textbf{[K]} In the T0 framework all observables can be formulated as ratios that
require no absolute mass parameter. The basic idea: instead of carrying $m$ in SI
units as a parameter, everything is expressed through $\xi$ and time intervals.

\textbf{[K]} \textbf{Systematic mass elimination.} From the T0 duality
$\tilde T\cdot m=1$ it follows that $m=1/\tilde T$. Substituting into the Dirac
equation:
\[
  \bigl[i\gamma^\mu\partial_\mu - 1/\tilde T(x,t)\bigr]\psi = 0.
\]
The mass parameter disappears from the equation; what remains is the time field
$\tilde T(x,t)$ as a dynamic quantity. This is the fully mass-free T0 formulation.

\textbf{[S]} \textbf{Philosophical consequence.} If all physics is formulable in
ratios (without absolute masses), then masses are not fundamental objects but
auxiliary constructs of measurement. What is fundamental: time intervals and the
geometry of the torus. This is the same thesis as A085 (natural units) in
relativistic language.

\section{Mass-Proportional Coupling}

\textbf{[K]} In the T0 Lagrangian the time field couples mass-proportionally:
$g_T^\ell = m_\ell\cdot\xi$. This is the mechanism behind the $g$-2 calculation
(A138): the coupling is not a free Yukawa constant but fully determined by the
mass and $\xi$. The one-loop formula for the anomalous magnetic moment follows
directly from this mass-proportional coupling.

\section{Ratios vs. Absolute Values in the Dirac Structure}

\textbf{[B]} The same distinction as in A261 and A138: ratios of Dirac corrections
(e.g.\ $a_\mu/a_e$) are correction-free and follow from the structure of the
Clifford algebra and the quantum numbers. Absolute values of the $g$-2 terms
require the SI bridge constants and $K_\text{frak}$ (A138).

\section{Verification Script}

\begin{itemize}[leftmargin=2em,itemsep=2pt,topsep=2pt]
  \item Numerical verification: standard Dirac equation as limiting case of the T0
    form ($\partial_\mu m\to0$); spin-$\tfrac12$ from winding $(1,1)$; mass-free
    formulation for simple cases:
    \texttt{python/A\_Serie\_Skripte/a263\_dirac.py}
\end{itemize}

\section{Open Questions}

\textbf{The complete mass-free formulation for all interactions is missing.} Mass
elimination is shown for free fields; for interaction terms it has not been
systematically carried out.

\textbf{The quark sector is missing.} The Dirac equation for quarks with dynamic
mass and colour charge is not developed at the same level as for leptons (A150).

\textbf{The relation to QED renormalisation is not worked out.} How the T0 Dirac
equation relates to perturbative QED (Feynman rules, loop integrals) is not
systematically developed (A075, open Hopf algebra).

\section{Supersedes}

This document subsumes: Dok.~051 (T0 Dirac equation, Clifford algebra, spin as
topology, mass-proportional coupling, ratios vs.\ absolute values), Dok.~052
(systematic mass elimination, complete mass-free T0 formulation), Dok.~053 (mass
elimination in the Dirac Lagrangian, ratio-based physics), Dok.~050 (simplified
Dirac equation).

\section{Use of AI Tools}

This document was created with the assistance of AI tools. Responsibility
for content and errors rests with the author. Full wording of the rules in A010.
